\documentclass[12pt, a4paper]{book}
\usepackage[margin=2.5cm]{geometry}
\usepackage{ctex}
\usepackage{amsmath, amssymb}
\usepackage{graphicx}
\usepackage{tabularx}
\usepackage{xcolor}
\usepackage{listings}
\usepackage{hyperref}
\usepackage{tikz}
\usetikzlibrary{shapes, arrows.meta, positioning, calc}

% listings 代码高亮配置
\lstset{
    language=Python,
    basicstyle=\ttfamily\small,
    keywordstyle=\color{blue},
    commentstyle=\color{green!50!black},
    stringstyle=\color{red},
    showstringspaces=false,
    breaklines=true,
    numbers=left,
    numberstyle=\tiny\color{gray},
    frame=single
}

% TikZ 样式定义（用于 Transformer 等架构图）
\tikzset{
    layer/.style={rectangle, draw, thick, minimum width=3cm, minimum height=0.8cm, align=center, rounded corners=2pt},
    attn/.style={layer, fill=blue!10},
    ffn/.style={layer, fill=orange!10},
    norm/.style={layer, fill=green!10},
    emb/.style={layer, fill=gray!20},
    pos/.style={circle, draw, thick, inner sep=2pt, fill=yellow!20},
    arrow/.style={-Stealth, thick}
}

\begin{document}
\tableofcontents
\newpage
\chapter{Activation\_functions}


\section{核心定义}
\textbf{中文定义}：激活函数是人工神经网络中负责将线性组合结果映射为非线性输出的数学函数。它决定了神经元是否被“激活”以及信号传递的程度，是神经网络具备模拟复杂非线性逻辑能力的关键。

\textbf{英文定义}：An activation function is a mathematical function that maps the linear combination of inputs to a non-linear output in an artificial neural network. It determines whether and to what extent a neuron should be "activated," enabling the network to learn complex non-linear patterns.

\textbf{为什么需要非线性？}：
根据\textbf{通用近似定理 (Universal Approximation Theorem)}，若没有非线性激活函数，无论神经网络有多少层，其最终输出始终是输入的线性组合（相当于单层感知机），无法拟合复杂的现实数据。

\section{核心原理与常用函数}

\subsection{1. Sigmoid 函数}
$$ \sigma(x) = \frac{1}{1 + e^{-x}} $$
\textbf{特点}：输出范围 $(0, 1)$，常用于二分类输出层。
\textbf{缺点}：存在\textbf{梯度消失}问题（当 $x$ 很大或很小时，梯度接近 0）；输出不以 0 为中心，导致收敛慢。

\subsection{2. Tanh 函数 (双曲正切)}
$$ \text{tanh}(x) = \frac{e^x - e^{-x}}{e^x + e^{-x}} $$
\textbf{特点}：输出范围 $(-1, 1)$，以 0 为中心，收敛通常比 Sigmoid 快。
\textbf{缺点}：依然存在梯度消失问题。

\subsection{3. ReLU (Rectified Linear Unit)}
$$ f(x) = \max(0, x) $$
\textbf{特点}：计算简单，有效缓解了正向区域的梯度消失，是目前 CNN 中的事实标准。
\textbf{缺点}：\textbf{Dead ReLU} 问题（导数为 0 的神经元可能永远无法更新）。

\subsection{4. Leaky ReLU}
$$ f(x) = \max(ax, x), \quad 0 < a < 1 $$
\textbf{特点}：解决了 Dead ReLU 问题，给负半轴分配一个微小的斜率。

\subsection{5. GELU (Gaussian Error Linear Unit)}
$$ f(x) = x \cdot \Phi(x) = x \cdot \frac{1}{2} [1 + \text{erf}(x/\sqrt{2})] $$
\textbf{特点}：引入了随机正则化思想，输出更平滑，是 Transformer (BERT, GPT) 的标配。

\section{严谨数学公式与导数推导}
在反向传播中，激活函数的导数决定了梯度传递的速度：

\begin{itemize}
    \item \textbf{Sigmoid 导数}：$\sigma'(x) = \sigma(x)(1 - \sigma(x))$。其最大值为 0.25。
    \item \textbf{ReLU 导数}：$\begin{cases} 1 & x > 0 \\ 0 & x < 0 \end{cases}$。
    \item \textbf{GELU 导数}：其导数不具有简单的闭式解，但其平滑特性使得模型在训练初期更稳定。
\end{itemize}

\section{模型结构中的位置}
激活函数通常紧跟在“权重加权和 (\text{Weight Sum})”之后：
\begin{verbatim}
[ 输入 X ] --> [ W * X + b ] --> [ 激活函数 f ] --> [ 输出 a ]
\end{verbatim}

\section{关键代码片段 (PyTorch实现)}
\begin{lstlisting}
import torch
import torch.nn as nn

# 1. 函数式接口
x = torch.randn(2, 2)
output_relu = torch.relu(x)

# 2. 模块化定义
model = nn.Sequential(
    nn.Linear(10, 20),
    nn.ReLU(),       # 基础非线性
    nn.Linear(20, 5),
    nn.LeakyReLU(0.01), # 解决 Dead ReLU
    nn.GELU(),       # Transformer 常用
    nn.Softmax(dim=-1)  # 多分类概率输出
)
\end{lstlisting}

\section{深度面试题与解答}
\subsection{基础概念}
\textbf{问：为什么非线性激活函数对深度学习至关重要？}\\
答：因为线性算子（矩阵乘法）具有结合律。多层线性层的叠加 $W_2(W_1 x)$ 等价于单个线性层 $W_{total} x$。没有激活函数，深度模型将退化为线性回归。

\subsection{进阶原理}
\textbf{问：ReLU 为什么能解决梯度消失问题？}\\
答：在 $x > 0$ 的区域，ReLU 的导数恒为 1。这意味着梯度在反向传播经过该神经元时不会被缩减，不仅计算极快，而且允许训练非常深的网络。

\subsection{工程实战}
\textbf{问：如何选择激活函数？}\\
答：经验法则是：
1. 隐藏层优先使用 \textbf{ReLU} 或其变体（Leaky ReLU/ELU）。
2. Transformer 架构优先选 \textbf{GELU}。
3. 二分类输出层用 \textbf{Sigmoid}，多分类用 \textbf{Softmax}。
4. 回归任务输出层通常不加激活函数。

\subsection{坑点分析}
\textbf{问：Softmax 和 Sigmoid 在损失计算上有什么配合细节？}\\
答：在 PyTorch 中，使用 \texttt{CrossEntropyLoss} 时，内部已经继承了 \texttt{LogSoftmax}。因此在定义网络模型时，最后一层输出层\textbf{不应}手动添加 Softmax，否则会重复计算导致性能下降或数值不稳定。

\section{优缺点对比表}

\begin{table}[h]
\centering
\begin{tabularx}{\textwidth}{|l|c|X|X|}
\hline
\textbf{函数} & \textbf{范围} & \textbf{主要优势} & \textbf{关键缺陷} \\
\hline
Sigmoid & (0, 1) & 概率解释性强 & 容易梯度消失，非零中心 \\
\hline
Tanh & (-1, 1) & 零中心输出 & 依然有梯度消失问题 \\
\hline
ReLU & [0, +$\infty$) & 计算极快，缓解梯度消失 & Dead ReLU (神经元“死亡”) \\
\hline
GELU & (-0.17, +$\infty$) & 性能优异，平滑度高 & 计算量略大于 ReLU \\
\hline
\end{tabularx}
\end{table}

\section{参考文献与拓展}
\begin{enumerate}
    \item \textbf{经典文献}: Glorot, X., et al. (2011). \textit{Deep Sparse Rectifier Neural Networks}. (ReLU 论文)
    \item \textbf{GELU 论文}: Hendrycks, D., & Gimpel, K. (2016). \textit{Gaussian Error Linear Units (GELUs)}.
    \item \textbf{CS231n}: Activation Functions and Their Performance Comparison (Stanford University).
\end{enumerate}

\chapter{Bert}


\section{核心定义}
\textbf{中文定义}：BERT（来自 Transformer 的双向编码器表征）是由 Google 在 2018 年提出的一种预训练语言模型。其核心架构是 Transformer 的 \textbf{Encoder-only} 结构，通过在超大规模语料上进行无监督预训练（MLM 和 NSP 任务），学习深层的双向上下文特征表示。

\textbf{英文定义}：BERT (Bidirectional Encoder Representations from Transformers) is a pre-trained NLP model introduced by Google. It uses the Encoder-only Transformer architecture and is trained on massive unlabeled text using Masked Language Modeling (MLM) and Next Sentence Prediction (NSP) to learn deep bidirectional representations.

\textbf{核心哲学}：
“双向性”是 BERT 的灵魂。不同于传统的自回归模型（如 GPT）只能从左向右看，BERT 同时利用左右两侧的上下文，极大地提升了模型对语言深层语义的理解能力。

\section{输入表示：三层 Embedding 叠加}
BERT 的输入不仅仅是词向量，而是三个层级的 Embedding 相加：
\begin{enumerate}
    \item \textbf{Token Embeddings}：WordPiece 分词后的向量。首个 Token 始终为 \texttt{[CLS]}（用于下游任务分类），句子间用 \texttt{[SEP]} 分隔。
    \item \textbf{Segment Embeddings}：区分两个不同句子的向量（A 句或 B 句），用于处理问答、内切等句对任务。
    \item \textbf{Position Embeddings}：学习到的位置向量，用于赋予模型位置感知能力。
\end{enumerate}

\section{两大预训练任务}

\subsection{1. MLM (Masked Language Model)}
\textbf{完形填空}：随机遮盖输入序列中 $15\%$ 的 Token，让模型根据上下文预测被遮盖的词。
\begin{itemize}
    \item $80\%$ 的概率替换为 \texttt{[MASK]}。
    \item $10\%$ 的概率替换为随机的一个词。
    \item $10\%$ 的概率保持原词。
\end{itemize}
\textbf{目的}：强制编码器学习双向的词义信息，解决由于双向扫描可能带来的“信息泄露”问题。

\subsection{2. NSP (Next Sentence Prediction)}
\textbf{下句预测}：给定句对 A 和 B，预测 B 是否是 A 的下一句。这是一个二分类任务。
\begin{itemize}
    \item $50\%$ 的样本 B 是 A 的真下句（Label: IsNext）。
    \item $50\%$ 的样本 B 是语料库中随机抽取的其他句子（Label: NotNext）。
\end{itemize}
\textbf{目的}：学习句子粒度之间的关系（Sentence-level coherence），对 QA、NLI 任务至关重要。

\section{微调范式 (Fine-tuning)}
BERT 的强大之处在于其“极简”的下游任务适配：
\begin{enumerate}
    \item \textbf{单句分类}：取首个 Token \texttt{[CLS]} 的输出向量，接全连接层。
    \item **标注任务**（如 NER）：取每个 Token 的输出向量进行序列预测。
    \item **句对匹配**：将两句拼在一起，取 \texttt{[CLS]} 输出。
    \item **问答任务**（如 SQuAD）：让两个全连接层分别预测答案在段落中的起始和结束位置。
\end{enumerate}

\section{BERT 家族主要变体}
\begin{itemize}
    \item \textbf{RoBERTa}：移除了 NSP 任务，使用更大的 Batch、更长的时间和更多的数据进行训练。
    \item \textbf{ALBERT}：通过参数共享和嵌入层分解极大减少了参数量，提升了训练速度。
    \item \textbf{DistilBERT}：利用知识蒸馏技术得到的轻量化版本，保留了 BERT $97\%$ 的性能但速度快了 $60\%$。
\end{itemize}

\section{关键代码：Masked 逻辑演示}
\begin{lstlisting}
import torch

# 假设输入 inputs_ids 为 [Batch, Seq]
# 创建遮盖矩阵 labels 记录真实 ID, inputs_ids 中被遮处填充 [MASK] 的 ID
mask_prob = 0.15
rand = torch.rand(input_ids.shape)
# 创建掩码
mask_arr = (rand < mask_prob) * (input_ids != 101) * (input_ids != 102) # 避开 [CLS], [SEP]

selection = torch.flatten((mask_arr).nonzero()).tolist()
input_ids[selection] = 103 # 103 是 [MASK] 的 ID
\end{lstlisting}

\section{BERT vs. GPT}
\begin{table}[h]
\centering
\begin{tabularx}{\textwidth}{|l|X|X|}
\hline
\textbf{特性} & \textbf{BERT} & \textbf{GPT} \\
\hline
\textbf{结构} & Encoder-only (Transformer) & Decoder-only (Transformer) \\
\hline
\textbf{上下文方向} & Bidirectional (双向) & Unidirectional (单向自回归) \\
\hline
\textbf{主要目标} & 理解、分类、特征提取 & 生成、创作、对话 \\
\hline
\textbf{推理效率} & 较高 (批量并行) & 逐词生成，KV-cache 优化 \\
\hline
\end{tabularx}
\end{table}

\section{深度面试题}

\textbf{问：BERT 为什么不用 Transformer 的 Decoder 结构？}\\
答：Decoder 结构的 Self-Attention 有掩码，只能看到左侧，这会限制模型对“全貌”的理解。BERT 的目标是提取最优表征，因此必须看到全文，Encoder 是天然的选择。

\textbf{问：RoBERTa 为什么要移除 NSP 任务？}\\
答：实验证明，对于长序列预训练，NSP 任务有时会损害单句表示的可学习性。RoBERTa 通过增大训练数据和步数，即便没有 NSP，其下游任务表现也全面超越了 BERT。

\section{参考文献}
\begin{enumerate}
    \item Devlin, J., et al. (2018). \textit{BERT: Pre-training of Deep Bidirectional Transformers for Language Understanding}.
    \item Liu, Y., et al. (2019). \textit{RoBERTa: A Robustly Optimized BERT Pretraining Approach}.
\end{enumerate}

\chapter{Cnn}


\section{核心定义}
\textbf{中文定义}：卷积神经网络（Convolutional Neural Network, CNN）是一种专门用于处理具有网格结构数据（如图像、视频）的深度神经网络。它通过卷积运算提取局部特征，并利用参数共享和池化操作降低模型复杂度。

\textbf{英文定义}：Convolutional Neural Networks (CNNs) are a class of deep neural networks specifically designed to process data with a grid-like topology (e.g., images). They extract spatial features through convolution operations and reduce complexity using parameter sharing and pooling.

\textbf{提出时间与作者}：现代 CNN 的基础是由 Yann LeCun 等人在 1998 年提出的 LeNet-5 架构确立的。其爆发点是 2012 年 AlexNet 在 ImageNet 比赛中的巨大成功。

\textbf{原始关键论文}：\textit{Gradient-Based Learning Applied to Document Recognition} (LeCun et al., 1998)

\textbf{技术体系定位及历史演进}：
CNN 起源于对生物视觉皮层的研究（感受野概念）。其演进历程经历了从早期的 LeNet 到深层的 VGG、残差连接的 ResNet、追求效率的 MobileNet，以及近期将注意力机制融入卷积的变体（如 ConvNeXt）。CNN 的核心价值在于将复杂的图像处理从繁琐的手工特征提取转向了自动化的多层表征学习。

\section{核心原理与底层逻辑}
CNN 的三大核心支柱是：\textbf{局部感受野（Local Receptive Fields）}、\textbf{权值共享（Shared Weights）}和\textbf{池化/下采样（Pooling）}。

\subsection{数据流向简述}
假定输入图像尺寸为 $(C, H, W)$（通道数、高度、宽度）。
\begin{enumerate}
    \item \textbf{卷积层}：使用多个大小为 $(C, k, k)$ 的卷积核在输入上滑动，执行点积运算得到特征图。
    \item \textbf{非线性激活}：通常接 ReLU 函数，增强模型的非线性表达能力。
    \item \textbf{池化层}：在特征图上进行最大值或平均值采样，减小空间维度 $(H, W)$，同时保持平移不变性。
    \item \textbf{全连接层}：经过多轮“卷积+池化”后，将多维张量展平（Flatten）为向量，输出分类或回归结果。
\end{enumerate}

\subsection{算法执行步序}
\begin{enumerate}
    \item 准备输入图像阵列。
    \item 应用卷积核执行互相关运算。
    \item 加上偏置项（Bias）并执行激活函数。
    \item 执行池化操作降低参数量。
    \item 重复上述步骤提取高层抽象特征。
    \item 最终特征送入全连接网络（MLP）进行预测。
\end{enumerate}

\section{严谨数学公式与推导}
二维卷积运算的核心公式为：
$$
Y(i,j) = (X * K)(i,j) = \sum_{m=0}^{k-1} \sum_{n=0}^{k-1} X(i+m, j+n) \cdot K(m,n)
$$
\textbf{符号说明}：
\begin{itemize}
    \item $X$：输入特征图。
    \item $K$：卷积核（Kernel），大小为 $k \times k$。
    \item $Y$：输出特征图中的一点。
\end{itemize}

\textbf{输出尺寸计算公式（至关重要）}：
若输入大小为 $W \times H$，卷积核为 $F \times F$，步长为 $S$，填充为 $P$，则输出宽度 $W_{out}$ 为：
$$
W_{out} = \left\lfloor \frac{W + 2P - F}{S} \right\rfloor + 1
$$
同理可得 $H_{out}$。

\section{模型结构/架构示意图}
典型的经典 CNN（如 LeNet 风格）架构如下：
\begin{verbatim}
[输入图像: 1x28x28]
      |
[ 5x5 Conv, 6 filters, Stride 1 ] --> (6x24x24)
      |
[ 2x2 MaxPool, Stride 2 ] ----------> (6x12x12)
      |
[ 5x5 Conv, 16 filters, Stride 1 ] -> (16x8x8)
      |
[ 2x2 MaxPool, Stride 2 ] ----------> (16x4x4)
      |
[ Flatten & Fully Connected ] ------> (120) -> (84)
      |
[ Output Classifier ] --------------> (10)
\end{verbatim}

\section{关键代码片段 (PyTorch实现)}
以下为使用 PyTorch 定义一个基础 CNN 模块的示例：

\begin{lstlisting}
import torch.nn as nn
import torch.nn.functional as F

class SimpleCNN(nn.Module):
    def __init__(self, num_classes=10):
        super(SimpleCNN, self).__init__()
        # 卷积层: 输入通道1, 输出通道16, 卷积核3x3, 填充1
        # 输入:(B, 1, 28, 28) -> 输出:(B, 16, 28, 28)
        self.conv1 = nn.Conv2d(1, 16, kernel_size=3, padding=1)
        
        # 池化层: 2x2 窗口, 步长2
        # 输入:(B, 16, 28, 28) -> 输出:(B, 16, 14, 14)
        self.pool = nn.MaxPool2d(2, 2)
        
        # 全连接层
        self.fc = nn.Linear(16 * 14 * 14, num_classes)

    def forward(self, x):
        # x: (B, 1, 28, 28)
        x = self.pool(F.relu(self.conv1(x))) # 卷积 -> 激活 -> 池化
        x = x.view(x.size(0), -1)           # 展平
        x = self.fc(x)                      # 分类
        return x
\end{lstlisting}

\section{深度面试题与解答}
\subsection{基础概念}
\textbf{问：CNN 相比普通的全连接神经网络有什么优势？}\\
答：主要有三点：1. \textbf{参数共享}（卷积核在全图移动使用同一组权重），极大减少了参数量；2. \textbf{局部连接}（每个神经元只看局部像素），符合视觉处理逻辑并提取空间特征；3. \textbf{平移不变性}（通过池化实现），图像中物体移动一小段距离不影响识别。

\subsection{进阶原理}
\textbf{问：1x1 卷积核的具体作用是什么？}\\
答：1x1 卷积不改变空间尺寸（H, W），但主要用于：1. \textbf{通道升降维}（控制计算负载，如 Inception 的 Bottleneck）；2. \textbf{跨通道信息融合}（让不同通道的特征进行组合）；3. \textbf{在不增加感受野的前提下增加非线性层}。

\subsection{工程实战}
\textbf{问：空洞卷积（Dilated Convolution）解决了什么问题？}\\
答：它通过在卷积核元素间插入“洞”来扩大感受野，关键优势是在\textbf{不增加参数量}且\textbf{不损失空间分辨率}（不使用下采样）的情况下，捕获更远距离的信息，常用于图像分割任务。

\subsection{坑点分析}
\textbf{问：池化层虽然好用，但为什么有些现代架构（如 ResNet）倾向于减少池化的使用？}\\
答：池化会带来不可逆的信息损失。现代常用“步长为 2 的卷积”代替池化来实现下采样。这样模型可以学习如何去下采样，而不是暴力丢弃（池化）像素，在保留细微空间特征上表现更好。

\section{优缺点与适用场景}

\begin{table}[h]
\centering
\begin{tabularx}{\textwidth}{|l|X|}
\hline
\textbf{维度} & \textbf{分析} \\
\hline
\textbf{优势} & 自动且高效的特征提取；强大的空间表征能力；模型参数量由于共享权重而高度精简。 \\
\hline
\textbf{局限} & 难以处理非网格数据（如社交网络图）；对物体的旋转、缩放敏感度较高（除非做大量增强）。 \\
\hline
\textbf{最佳适用场景} & CV 领域几乎所有任务（识别、检测、分割）；音频频谱处理；视频分析。 \\
\hline
\textbf{替代方案} & 极致长跨度依赖任务使用 Transformer；非结构化图数据使用 GNN。 \\
\hline
\end{tabularx}
\end{table}

\section{参考文献与拓展}
\begin{enumerate}
    \item \textbf{里程碑论文}: Krizhevsky, A., et al. (2012). \textit{ImageNet classification with deep convolutional neural networks}. (AlexNet)
    \item \textbf{必备课程}: CS231n: Convolutional Neural Networks for Visual Recognition (Stanford).
    \item \textbf{变体探索}: He, K., et al. (2016). \textit{Deep Residual Learning for Image Recognition}. (ResNet)
\end{enumerate}

\chapter{Diffusion}


\section{核心定义}
\textbf{中文定义}：扩散模型（Diffusion Models）是一类生成模型，其灵感源自非平衡热力学。它通过定义一个正向扩散过程（逐渐向数据中添加高斯噪声直至其变为纯噪声）和一个反向去噪过程（学习从噪声中恢复原始数据）来学习复杂的数据分布。

\textbf{英文定义}：Diffusion Models are a class of generative models inspired by non-equilibrium thermodynamics. They work by defining a forward diffusion process that systematically adds Gaussian noise to data and a learned reverse process that restores the data from noise.

\textbf{代表作}：DDPM (Denoising Diffusion Probabilistic Models, 2020)。

\section{核心原理：前向与反向过程}
扩散模型的核心由两个互逆的过程组成：

\subsection{1. 前向过程 (Forward Diffusion)}
将原始数据 $x_0$ 逐步转变为纯高斯噪声 $x_T$。每一步添加微小噪声：
$$ q(x_t | x_{t-1}) = \mathcal{N}(x_t; \sqrt{1 - \beta_t} x_{t-1}, \beta_t \mathbf{I}) $$
利用重参数化技巧，可以直接从 $x_0$ 计算任意时刻的 $x_t$：
$$ x_t = \sqrt{\bar{\alpha}_t} x_0 + \sqrt{1 - \bar{\alpha}_t} \epsilon, \quad \epsilon \sim \mathcal{N}(0, \mathbf{I}) $$
其中 $\alpha_t = 1 - \beta_t$，$\bar{\alpha}_t = \prod_{i=1}^t \alpha_i$。

\subsection{2. 反向过程 (Reverse Denoising)}
模型学习预测噪声 $\epsilon_\theta(x_t, t)$，从而实现从 $x_t$ 到 $x_{t-1}$ 的逐步还原：
$$ p_\theta(x_{t-1} | x_t) = \mathcal{N}(x_{t-1}; \mu_\theta(x_t, t), \Sigma_\theta(x_t, t)) $$

\section{训练目标函数}
DDPM 的简化损失函数通常采用均方误差（MSE），目标是让网络预测注入的噪声 $\epsilon$：
$$ L_{simple}(\theta) = \mathbb{E}_{x_0, \epsilon, t} \left[ \| \epsilon - \epsilon_\theta(\sqrt{\bar{\alpha}_t} x_0 + \sqrt{1 - \bar{\alpha}_t} \epsilon, t) \|^2 \right] $$

\section{模型结构：时间敏感型 U-Net}
扩散模型通常使用 U-Net 架构，关键改进是引入了\textbf{时间嵌入 (Time Embedding)}，使模型知道当前处于哪个去噪阶段。
\begin{verbatim}
[ 噪声图像 x_t ] --+--> [ Encoder (下采样) ] --+--> [ Bottleneck ]
                    |           |               |
[ 时间步 t ] -------+--> [ 加入 Time Embedding ] +--> [ Decoder (上采样) ]
                                                |
                                                v
                                         [ 预测的噪声 epsilon ]
\end{verbatim}

\section{关键代码片段 (PyTorch简化实现)}
\begin{lstlisting}
import torch
import torch.nn as nn

class DiffusionTrainer:
    def __init__(self, model, betas):
        self.model = model
        self.alphas = 1.0 - betas
        self.alphas_cumprod = torch.cumprod(self.alphas, dim=0)

    def get_loss(self, x_0):
        batch_size = x_0.shape[0]
        # 1. 随机采样时间步 t
        t = torch.randint(0, len(self.alphas), (batch_size,))
        # 2. 采样标准正态噪声
        noise = torch.randn_like(x_0)
        
        # 3. 计算含噪图像 x_t
        sqrt_alphas_cumprod_t = self.alphas_cumprod[t].sqrt().view(-1, 1, 1, 1)
        sqrt_one_minus_alphas_cumprod_t = (1 - self.alphas_cumprod[t]).sqrt().view(-1, 1, 1, 1)
        x_t = sqrt_alphas_cumprod_t * x_0 + sqrt_one_minus_alphas_cumprod_t * noise
        
        # 4. 模型预测噪声并计算 MSE Loss
        predicted_noise = self.model(x_t, t)
        return torch.nn.functional.mse_loss(noise, predicted_noise)
\end{lstlisting}

\section{深度面试题与解答}
\subsection{基础概念}
\textbf{问：扩散模型相比 GAN 有什么核心优势？}\\
答：1. \textbf{训练极度稳定}：Diffusion 是基于似然概率的显式建模，不存在 GAN 那种难平衡的对抗博弈和模式崩溃。2. \textbf{生成多样性高}：Diffusion 能更好地覆盖数据分布的所有模态。

\subsection{进阶原理}
\textbf{问：什么是 Classifier-Free Guidance (CFG)？为什么它很重要？}\\
答：它是一种在推理时不需要额外分类器就能大幅提升生成质量（文本一致性）的技术。通过同时训练有条件和无条件的模型（即将条件设为空），并在生成时将两者的预测方向进行外推：$\tilde{\epsilon} = \epsilon_{uncond} + w \cdot (\epsilon_{cond} - \epsilon_{uncond})$。

\subsection{工程实战}
\textbf{问：LDM (Latent Diffusion Model) 的贡献是什么？}\\
答：它将扩散过程从像素空间放到了预训练 VAE 的\textbf{潜空间 (Latent Space)}。这极大地降低了计算开销，使得在消费级 GPU 上进行大规模高分辨率图像生成成为可能（Stable Diffusion 的基础）。

\subsection{坑点分析}
\textbf{问：Diffusion 推理速度慢的问题如何解决？}\\
答：原始 DDPM 需要几百上千步迭代。解决方法包括：1. \textbf{DDIM}（确定性采样，仅需 20-50 步）；2. \textbf{蒸馏技术 (Distillation)}，将多步采样能力通过教师模型教给只需几步的弟子模型。

\section{生成模型横向对比表}

\begin{table}[h]
\centering
\begin{tabularx}{\textwidth}{|l|X|X|X|}
\hline
\textbf{维度} & \textbf{VAE} & \textbf{GAN} & \textbf{Diffusion} \\
\hline
训练目标 & 变分下界最大化 & 对抗博弈 & 噪声预测 (MSE) \\
\hline
样本质量 & 模糊 (Over-smoothed) & 高 (可能缺乏多样性) & 极高且多样 \\
\hline
生成速度 & 快 (单次前向) & 快 (单次前向) & 慢 (多次迭代) \\
\hline
稳定性 & 稳定 & 不稳定 & 非常稳定 \\
\hline
\end{tabularx}
\end{table}

\section{参考文献与拓展}
\begin{enumerate}
    \item \textbf{里程碑论文}: Ho, J., et al. (2020). \textit{Denoising Diffusion Probabilistic Models}.
    \item \textbf{Latent Diffusion}: Rombach, R., et al. (2022). \textit{High-Resolution Image Synthesis with Latent Diffusion Models}.
    \item \textbf{进阶资源}: \textit{What are Diffusion Models?} (Lilian Weng's Blog, 深度解析数学原理)。
\end{enumerate}

\chapter{Gan}


\section{核心定义}
\textbf{中文定义}：生成对抗网络（Generative Adversarial Network, GAN）是一种通过对抗训练学习数据分布的深度学习架构。它由生成器（Generator）和判别器（Discriminator）组成，两者在相互博弈中共同进化：生成器学习制造逼真数据，判别器学习分辨真伪。

\textbf{英文定义}：Generative Adversarial Networks (GANs) are a framework for training generative models based on a game-theoretic scenario. It consists of two networks: a Generator that learns to produce realistic data and a Discriminator that learns to distinguish between real and generated samples.

\textbf{提出时间与作者}：2014年，由 Ian Goodfellow 等人提出。

\textbf{原始关键论文}：\textit{Generative Adversarial Nets} (NIPS 2014)

\section{核心原理：零和博弈}
GAN 的核心在于其独特的目标函数，定义为一个\textbf{极小极大博弈 (Minimax Game)}。

\subsection{1. 目标函数}
$$ \min_G \max_D V(D, G) = \mathbb{E}_{x \sim p_{data}(x)}[\log D(x)] + \mathbb{E}_{z \sim p_z(z)}[\log(1 - D(G(z)))] $$
\textbf{符号说明}：
\begin{itemize}
    \item $D(x)$：判别器，输出 $x$ 来自真实数据的概率。
    \item $G(z)$：生成器，将随机噪声 $z$ 映射到数据空间。
    \item $\max_D$：判别器希望最大化准确率（真实样本输出近 1，生成样本输出近 0）。
    \item $\min_G$：生成器希望最小化判别器的判断力（使得生成样本让判别器输出近 1）。
\end{itemize}

\subsection{2. 训练过程}
训练是一个交替进行的迭代过程：
\begin{enumerate}
    \item \textbf{固定 G，训练 D}：使判别器更强。
    \item \textbf{固定 D，训练 G}：使生成器造出的图更像真的，从而“欺骗”判别器。
    \item \textbf{纳什均衡}：理论上当 $p_g = p_{data}$ 时，判别器输出恒为 0.5，此时达到平衡。
\end{enumerate}

\section{模型结构示意图}
\begin{verbatim}
[ 随机噪声 z ] --> [ 生成器 G ] --> [ 生成样本 G(z) ] 
                                          |
                                          v
      [ 真实样本 x ] ---------------> [ 判别器 D ] --> [ 分类: 真/假 ]
\end{verbatim}

\section{关键代码片段 (PyTorch实现)}
以下为基础 GAN 的核心训练逻辑简化版：
\begin{lstlisting}
import torch
import torch.nn as nn

# 定义损失函数
criterion = nn.BCELoss()

# 1. 训练判别器 D
# 真实数据 Loss
outputs = D(real_images)
d_loss_real = criterion(outputs, torch.ones_like(outputs))
# 生成数据 Loss
z = torch.randn(batch_size, latent_dim)
fake_images = G(z)
outputs = D(fake_images.detach()) # detach防止更新G
d_loss_fake = criterion(outputs, torch.zeros_like(outputs))
# 反向传播 D
d_loss = d_loss_real + d_loss_fake
d_optimizer.step()

# 2. 训练生成器 G
z = torch.randn(batch_size, latent_dim)
fake_images = G(z)
outputs = D(fake_images) # 重新通过最新的D
# 目标：让 D 认为 fake 是 1 (真)
g_loss = criterion(outputs, torch.ones_like(outputs))
# 反向传播 G
g_optimizer.step()
\end{lstlisting}

\section{深度面试题与解答}
\subsection{基础概念}
\textbf{问：为什么 GAN 的对抗博弈比普通网络更难训练？}\\
答：因为 GAN 不是寻找损失函数的全局最小值，而是寻找两者的纳什均衡点点。如果 D 太强，G 的梯度会消失；如果 D 太弱，G 无法学到任何有用的分布。

\subsection{进阶原理}
\textbf{问：什么是模式崩溃 (Mode Collapse)？}\\
答：指生成器由于发现判别器的某个漏洞，倾向于生成极少数甚至单一的、能够成功欺骗判别器的样本，而失去了数据分布的多样性（例如在 MNIST 上只生成数字 1）。

\subsection{工程实战}
\textbf{问：WGAN (Wasserstein GAN) 解决了什么核心问题？}\\
答：解决了原生 GAN 使用 JS 散度在分布不重叠时梯度消失的问题。WGAN 使用 \textbf{EM距离 (Earth-Mover's distance)} 配合 \textbf{Lipschitz 连续性约束 (Weight Clipping)}，提供了更平滑的梯度映射。

\subsection{坑点分析}
\textbf{问：在评估 GAN 生成质量时，为什么不用训练集 Loss？}\\
答：GAN 的 Loss 只能反映 D 与 G 的博弈进度，不能代表图像质量。常用 \textbf{FID (Fréchet Inception Distance)} 或 \textbf{Inception Score (IS)} 来定量评估图像的逼真度和多样性。

\section{GAN 变体对比表}

\begin{table}[h]
\centering
\begin{tabularx}{\textwidth}{|l|X|}
\hline
\textbf{变体} & \textbf{核心改进} \\
\hline
\textbf{DCGAN} & 引入卷积神经网络（CNN）与 Batch Normalization，确立了工程基础。 \\
\hline
\textbf{WGAN / WGAN-GP} & 引入 Wasserstein 距离与梯度惩罚，极大提升了训练稳定性。 \\
\hline
\textbf{CycleGAN} & 实现了无需配对数据的跨域图像风格转换。 \\
\hline
\textbf{StyleGAN} & 引入解耦控制空间，实现了极高分辨率、可控细节的人脸生成。 \\
\hline
\end{tabularx}
\end{table}

\section{参考文献与拓展}
\begin{enumerate}
    \item \textbf{奠基论文}: Goodfellow, I., et al. (2014). \textit{Generative Adversarial Nets}. NIPS.
    \item \textbf{工程标杆}: Radford, A., et al. (2015). \textit{Unsupervised Representation Learning with DCGAN}.
    \item \textit{GAN Lab}: Google 提供的交互式 GAN 训练可视化工具。
\end{enumerate}

\chapter{Gpt}


\section{核心定义}
\textbf{中文定义}：GPT（生成式预训练 Transformer）是由 OpenAI 提出的一系列自回归语言模型。其核心架构基于 Transformer 的 \textbf{Decoder-only} 结构，通过在大规模无标注文本上进行“下一个词预测”（Next Token Prediction）的生成式预训练，学习通用的语言表征能力。

\textbf{英文定义}：GPT (Generative Pre-trained Transformer) is a series of auto-regressive language models developed by OpenAI. Based on the Decoder-only Transformer architecture, it is pre-trained on massive unlabeled text via the objective of "next token prediction" to learn general-purpose language representations.

\textbf{核心哲学}：
GPT 的成功证明了“预训练 + 微调/零样本提示”范式的强大。它通过增大参数规模（Scaling Laws）和数据量，实现了从专用模型到通用人工智能（AGI）的跨越。

\section{架构：Decoder-only Transformer}

\subsection{为什么只用 Decoder？}
与原始 Transformer（Encoder-Decoder）和 BERT（Encoder-only）不同，GPT 舍弃了 Encoder，只保留 Decoder。
\begin{enumerate}
    \item \textbf{生成任务的天性}：生成任务是自左向右的，Decoder 的 Masked Self-Attention 完美契合这种训练方式。
    \item \textbf{自回归建模}：GPT 将条件概率分解为 $P(x) = \prod P(x_i | x_{<i})$。这种因果建模使得模型在推理时能够逐词生成。
\end{enumerate}

\subsection{结构组件}
一个典型的 GPT Block 包含：
\begin{itemize}
    \item \textbf{Masked Multi-Head Attention}：确保当前计算只能看到之前的词，防止“偷看答案”。
    \item \textbf{Layer Normalization}：通常采用 Pre-LN 结构以增加深层模型的稳定性。
    \item \textbf{Position-wise FFN}：两层全连接网络，通常使用 GeLU 激活函数。
\end{itemize}

\section{Evolution: 从 GPT-1 到 GPT-4}

\begin{table}[h]
\centering
\begin{tabularx}{\textwidth}{|l|l|X|}
\hline
\textbf{模型} & \textbf{参数量} & \textbf{核心贡献} \\
\hline
GPT-1 & 1.17 亿 & 确立了“生成式预训练 + 判别式微调”的二阶段流程。 \\
\hline
GPT-2 & 15 亿 & 强调 \textbf{Zero-shot} 能力，展示了模型在不经微调下处理多任务的潜力。 \\
\hline
GPT-3 & 1750 亿 & 引入 \textbf{In-Context Learning (ICL)}，通过 Few-shot 提示即可完成复杂任务。 \\
\hline
InstructGPT & 1750 亿 & 引入 \textbf{RLHF}，使模型输出符合人类指令和偏好（Alignment）。 \\
\hline
GPT-4 & 未公开 (万亿级) & 多模态理解力、处理长输入、极强的推理和逻辑能力。 \\
\hline
\end{tabularx}
\end{table}

\section{RLHF：对齐人类偏好}
为了让模型不产生胡言乱语或有害内容，OpenAI 引入了\textbf{人类反馈强化学习（RLHF）}：
\begin{enumerate}
    \item \textbf{监督微调 (SFT)}：在人工编写的问答对上进行微调。
    \item \textbf{奖励模型 (RM) 训练}：人类对模型的多个输出进行排序，训练一个评分模型。
    \item \textbf{强化学习 (PPO)}：使用奖励模型作为反馈，通过强化学习优化策略，使得生成内容得分最高。
\end{enumerate}

\section{GPT vs. BERT}
\begin{table}[h]
\centering
\begin{tabularx}{\textwidth}{|l|X|X|}
\hline
\textbf{特性} & \textbf{BERT} & \textbf{GPT} \\
\hline
\textbf{架构} & Encoder-only & Decoder-only \\
\hline
\textbf{注意力} & Bidirectional (双向) & Unidirectional (因果/单向) \\
\hline
\textbf{预训练目标} & MLM (填空) & CLM (预测下一个词) \\
\hline
\textbf{擅长领域} & 理解、分类、抽取 & 创意写作、对话、代码生成 \\
\hline
\end{tabularx}
\end{table}

\section{关键代码：Causal Self-Attention mask}
GPT 的核心在于其因果掩码，在 PyTorch 中通常这样实现：
\begin{lstlisting}
import torch
import torch.nn as nn

def create_causal_mask(size):
    # 生成一个上三角矩阵，并取反，使得左下角(包含对角线)为 1
    mask = torch.tril(torch.ones(size, size))
    return mask == 0 # True 的地方将被填充为 -inf

# 在 Attention 计算中的应用
# scores: [Batch, Heads, Seq, Seq]
mask = create_causal_mask(seq_len)
scores = scores.masked_fill(mask, float('-inf'))
attn = torch.softmax(scores, dim=-1)
\end{lstlisting}

\section{深度面试题}

\textbf{问：为什么 GPT 相比 BERT 展现出了更强的通用泛化能力？}\\
答：1. \textbf{预训练难度}：CLM（预测未来）比 MLM（填补中间）更难，模型被迫学习更深层的逻辑和语义；2. \textbf{任务统一}：自然语言的所有任务都可以转化为生成任务（即 Prompting），这使得 GPT 可以直接进行 Zero-shot 迁移，而 BERT 往往需要针对每个下游任务修改 Head 并全量微调。

\textbf{问：什么是涌现能力 (Emergent Abilities)？}\\
答：指当模型参数量突破某个临界点（如 600 亿）时，模型在某些任务（如逻辑推理、数学计算）上的表现会突然从接近随机猜测跃升到远超标准水平。这是大模型（LLM）区别于小模型的最显著特征。

\section{参考文献}
\begin{enumerate}
    \item Radford, A., et al. (2018). \textit{Improving Language Understanding by Generative Pre-Training} (GPT-1).
    \item Brown, T., et al. (2020). \textit{Language Models are Few-Shot Learners} (GPT-3).
    \item Ouyang, L., et al. (2022). \textit{Training language models to follow instructions with human feedback} (InstructGPT).
\end{enumerate}

\chapter{Information\_theory}


\section{核心定义}
\textbf{中文定义}：信息论是一门利用数学统计方法研究信息的量化、存储和通信的学科。在深度学习中，它提供了衡量概率分布差异（如 KL 散度）、评估模型不确定性（如信息熵）及优化目标函数（如交叉熵）的数学骨架。

\textbf{英文定义}：Information Theory is the mathematical study of the quantification, storage, and communication of information. In deep learning, it provides the framework for measuring differences between probability distributions (e.g., KL divergence) and optimizing objective functions.

\section{核心衡量指标}

\subsection{1. 自信息 (Self-Information)}
衡量一个事件发生的惊讶程度。事件 $x$ 的自信息为：
$$ I(x) = -\log p(x) $$
概率越小的事件，包含的信息量越大。

\subsection{2. 信息熵 (Entropy)}
所有可能事件自信息的期望值，衡量随机变量的不确定性：
$$ H(X) = \mathbb{E}_{x \sim P} [I(x)] = -\sum_{x \in \mathcal{X}} p(x) \log p(x) $$

\subsection{3. 相对熵 / KL 散度 (Relative Entropy / KL Divergence)}
衡量两个分布 $P$ 和 $Q$ 之间的差异：
$$ D_{KL}(P || Q) = \mathbb{E}_{x \sim P} \left[ \log \frac{p(x)}{q(x)} \right] = \sum p(x) \log \frac{p(x)}{q(x)} $$
\textbf{注意}：KL 散度是非对称的（即 $D_{KL}(P||Q) \neq D_{KL}(Q||P)$），且始终 $\geq 0$。

\subsection{4. 交叉熵 (Cross-Entropy)}
在深度学习分类任务中最为常用：
$$ H(P, Q) = -\mathbb{E}_{x \sim P} [\log q(x)] = -\sum p(x) \log q(x) $$
它可以拆解为：$H(P, Q) = H(P) + D_{KL}(P || Q)$。
当 $P$（真实标签）固定时，最小化交叉熵等价于最小化 $P, Q$ 之间的 KL 散度。

\subsection{5. 互信息 (Mutual Information)}
衡量两个变量之间的相互依赖程度：
$$ I(X; Y) = H(X) - H(X|Y) = H(Y) - H(Y|X) $$

\section{信息论与深度学习的关系}
\begin{itemize}
    \item \textbf{最大似然估计 (MLE)}：在监督学习中，最小化交叉熵等价于最大化似然函数。
    \item \textbf{VAE 优化}：VAE 的目标函数中包含 KL 散度项，用于约束隐空间分布。
    \item \textbf{信息瓶颈 (Information Bottleneck)}：一种解释深度网络学习过程的理论，认为网络在压缩输入信息的同时保留与输出最相关的特征。
\end{itemize}

\section{概念逻辑关系图}
\begin{verbatim}
[ 联合熵 H(X,Y) ]
      /      \
[ H(X) ] -- [ I(X;Y) ] -- [ H(Y) ]
      \      /
     [ H(X|Y) ]
\end{verbatim}

\section{代码片段 (熵与 KL 散度计算)}
\begin{lstlisting}
import torch
import torch.nn.functional as F

# 1. 计算交叉熵
logits = torch.tensor([[1.0, 2.0, 0.5]])
target = torch.tensor([1]) # 真实类别的索引
loss_ce = F.cross_entropy(logits, target)

# 2. 手动计算 KL 散度
p = torch.tensor([0.1, 0.9])
q = torch.tensor([0.2, 0.8])
kl_div = torch.sum(p * (torch.log(p) - torch.log(q)))

# 3. 使用 PyTorch 内置函数 (注意输入需要 log 空间)
kl_loss = F.kl_div(q.log(), p, reduction='batchmean')
\end{lstlisting}

\section{深度面试题与解答}
\subsection{基础概念}
\textbf{问：为什么 KL 散度不能被称为“距离”？}\\
答：因为“距离”必须满足\textbf{对称性}和\textbf{三角不等式}。KL 散度不满足对称性（$P$ 到 $Q$ 的差异不等于 $Q$ 到 $P$），也不满足三角不等式。它更准确的称呼是“信息增益”或“相对熵”。

\subsection{进阶原理}
\textbf{问：为什么在 GAN 中使用 JS 散度而不是 KL 散度？}\\
答：原生的 GAN 论文发现当 $P$ 和 $Q$ 几乎不重叠时，KL 散度是常数，导致梯度反传失败。JS 散度是 KL 散度的对称变体，取值在 $[0, \log 2]$ 之间，在分布不重叠时依然能提供更好的数值特性。

\subsection{工程实战}
\textbf{问：Label Smoothing（标签平滑）是如何从信息论角度解释的？}\\
答：标签平滑是将 One-hot 的 $P$ 变成一个更加熵增的分布。从交叉熵角度看，它限制了模型对样本的“极端自信”，增加了输出分布的熵值，从而起到正则化作用，防止过拟合。

\section{指标对比总结表}

\begin{table}[h]
\centering
\begin{tabularx}{\textwidth}{|l|X|l|}
\hline
\textbf{指标} & \textbf{核心用途} & \textbf{取值范围} \\
\hline
信息熵 & 衡量单一变量的不确定性 & $[0, \infty)$ \\
\hline
交叉熵 & 监督学习的分类损失函数 & $[0, \infty)$ \\
\hline
KL 散度 & 衡量分布间差异（生成模型常用） & $[0, \infty)$ \\
\hline
互信息 & 特征选择、表征学习相关性衡量 & $[0, \min(H(X),H(Y))]$ \\
\hline
\end{tabularx}
\end{table}

\section{参考文献与拓展}
\begin{enumerate}
    \item \textbf{奠基之作}: Shannon, C. E. (1948). \textit{A Mathematical Theory of Communication}.
    \item \textbf{经典教科书}: Cover, T. M., & Thomas, J. A. \textit{Elements of Information Theory}.
    \item \textbf{DL 应用}: Goodfellow, I., et al. \textit{Deep Learning} (Chapter 3: Probability and Information Theory).
\end{enumerate}

\chapter{Loss\_functions}


\section{核心定义}
\textbf{中文定义}：损失函数（Loss Function）是用于衡量模型预测值与真实标签之间差异的函数。它是深度学习优化过程的“指挥棒”，通过反向传播计算梯度，引导模型参数向减少误差的方向更新。

\textbf{英文定义}：A Loss Function (or Objective Function) is a mathematical function that quantifies the difference between the predicted value and the ground truth. It serves as the objective for optimization, guiding the model parameters to minimize errors through gradient descent.

\section{常用损失函数详解}

\subsection{1. 均方误差 (MSE)}
$$ L = \frac{1}{n} \sum_{i=1}^n (y_i - \hat{y}_i)^2 $$
\textbf{适用场景}：回归任务（如预测房价、坐标回归）。
\textbf{特点}：对异常值（Outliers）非常敏感。

\subsection{2. 二元交叉熵 (Binary Cross Entropy)}
$$ L = -\frac{1}{n} \sum_{i=1}^n [y_i \log(\hat{y}_i) + (1 - y_i) \log(1 - \hat{y}_i)] $$
\textbf{适用场景}：二分类任务。要求输出 $\hat{y}_i \in [0, 1]$（通常配对 Sigmoid）。

\subsection{3. 交叉熵 (Categorical Cross Entropy)}
$$ L = -\sum_{c=1}^M y_{o,c} \log(p_{o,c}) $$
\textbf{适用场景}：多分类任务。常与 Softmax 配合使用，衡量两个概率分布之间的距离。

\subsection{4. KL 散度 (Kullback-Leibler Divergence)}
$$ D_{KL}(P || Q) = \sum P(x) \log \frac{P(x)}{Q(x)} $$
\textbf{适用场景}：衡量两个概率分布的差异。在 VAE 中用于规范化隐空间分布。

\subsection{5. Smooth L1 损失 (Huber Loss)}
当误差小时呈 L2（平方），误差大时呈 L1（线性）。
\textbf{优点}：缓解了 MSE 对异常值的敏感，同时避免了 L1 在 0 点处不可导的问题。常用于目标检测中的边界框回归。

\section{反向传播中的角色}
参数更新公式为：$ W \leftarrow W - \eta \frac{\partial L}{\partial W} $。
损失函数的导数直接决定了初次回传的梯度大小。如果 Loss 选错（例如二分类用 MSE），会导致函数饱和区梯度极小，模型无法收敛。

\section{典型架构示意图}
\begin{verbatim}
[ 特征提取网络 ] --> [ 预测值 y_hat ] 
                            |
                            v
[ 真实标签 y ] ------> [ 损失函数 L ] --> [ 计算梯度 dL/dy_hat ]
                            |                   |
                            +-------------------+
                                     |
                                     v
                            [ 反向传播更新参数 ]
\end{verbatim}

\section{关键代码片段 (PyTorch实现)}
\begin{lstlisting}
import torch
import torch.nn as nn

# 1. 回归常用
mse_loss = nn.MSELoss()

# 2. 分类常用 (注意: CrossEntropyLoss 内部集成 LogSoftmax)
# 预测值应为 Raw Logits
ce_loss = nn.CrossEntropyLoss()

# 3. 二分类
# 如果输出已通过 Sigmoid:
bce_loss = nn.BCELoss()
# 如果输出是 Logits (推荐, 更数值稳定):
bce_with_logits = nn.BCEWithLogitsLoss()

# 使用示例
logits = model(input)
loss = ce_loss(logits, target)
loss.backward()
\end{lstlisting}

\section{深度面试题与解答}
\subsection{基础概念}
\textbf{问：分类任务为什么通常用交叉熵而不是 MSE？}\\
答：主要有两点：1. \textbf{学习效率}：交叉熵配合 Softmax 在反向传播时，梯度与预测误差成线性比例，即便误差很大也有强梯度；而 MSE 在分类问题中，输出层容易陷入饱和区，导致梯度消失。2. \textbf{概率意义}：交叉熵源于最大似然估计，更符合分布拟合的本质。

\subsection{进阶原理}
\textbf{问：Softmax + Cross-Entropy 的导数有什么优雅的形式？}\\
答：设预测概率为 $p_i$，真实标签为 $y_i$（One-hot），其梯度简写为 $\frac{\partial L}{\partial z_i} = p_i - y_i$。即预测值与真实值的差，简洁且直观。

\subsection{工程实战}
\textbf{问：样本不均衡（Imbalanced Data）时该如何调整损失函数？}\\
答：可以给不同类别设置权重（Weight），如 \texttt{nn.CrossEntropyLoss(weight=tensor)}。或者使用 \textbf{Focal Loss}，通过调整系数让模型更关注难分类的样本（Hard Examples）。

\subsection{坑点分析}
\textbf{问：L1 Loss 和 L2 Loss 分别有什么优缺点？}\\
答：L2 (MSE) 下梯度随误差减小而减小，有利于收敛，但易受离群点干扰。L1 梯度恒定，对异常值稳健，但在 0 点不可导且在小误差下容易震荡。

\section{损失函数对比总结表}

\begin{table}[h]
\centering
\begin{tabularx}{\textwidth}{|l|l|X|}
\hline
\textbf{函数} & \textbf{适用任务} & \textbf{主要优势} \\
\hline
MSE & 回归 & 凸函数，收敛平滑。 \\
\hline
Cross-Entropy & 分类 & 梯度传导快，信息论支持。 \\
\hline
Smooth L1 & 目标检测 & 稳健且在 0 附近平滑。 \\
\hline
Dice Loss & 语义分割 & 能够应对严重的类别不平衡。 \\
\hline
\end{tabularx}
\end{table}

\section{参考文献与拓展}
\begin{enumerate}
    \item \textbf{深度研究}: \textit{Deep Learning} (Ian Goodfellow) - Chapter 6.
    \item \textbf{PyTorch 文档}: \texttt{torch.nn} 模块下的 Loss Functions 章节。
    \item \textbf{进阶变体}: Lin, T. Y., et al. (2017). \textit{Focal Loss for Dense Object Detection}.
\end{enumerate}

\chapter{Mlp}


\section{核心定义}
\textbf{中文定义}：多层感知机（Multi-Layer Perceptron, MLP），又称人工神经网络（ANN）或全连接网络，是一种前馈神经网络。它由输入层、一个或多个隐藏层以及输出层组成，层与层之间通过全连接（Fully Connected）的方式连接，即上一层的所有神经元都与下一层的所有神经元相连。

\textbf{英文定义}：A Multi-Layer Perceptron (MLP) is a class of feedforward artificial neural network (ANN). It consists of at least three layers of nodes: an input layer, a hidden layer, and an output layer. Except for the input nodes, each node is a neuron that uses a nonlinear activation function.

\textbf{历史地位}：MLP 是现代深度学习的基石。1980年代反向传播（Backpropagation）算法的引入，使得训练深层 MLP 成为可能，从而克服了早期感知机无法解决非线性（如 XOR 问题）的局限。

\section{核心原理与计算过程}

\subsection{1. 层级结构}
\begin{itemize}
    \item \textbf{输入层}：接收原始特征向量 $x \in \mathbb{R}^d$。
    \item \textbf{隐藏层}：负责提取非线性特征。每个神经元执行：$a = \sigma(Wx + b)$。
    \item \textbf{输出层}：根据任务输出结果（如回归用线性，分类用 Softmax）。
\end{itemize}

\subsection{2. 计算公式}
对于第 $l$ 层，其输出 $h^{(l)}$ 的计算公式为：
$$ z^{(l)} = W^{(l)} h^{(l-1)} + b^{(l)} $$
$$ h^{(l)} = \sigma(z^{(l)}) $$
其中 $W^{(l)}$ 是权重矩阵，$b^{(l)}$ 是偏置项，$\sigma$ 是非线性激活函数。

\subsection{3. 通用近似定理 (Universal Approximation Theorem)}
该定理指出：一个拥有至少一个隐藏层（且包含足够多神经元）的前馈神经网络，只要配以合适的非线性激活函数，就可以以任意精度逼近闭区间上的任何连续函数。

\section{模型结构示意图}
\begin{verbatim}
[ Input x1 ] --\ /-- [ Hidden h1 ] --\ /-- [ Output y1 ]
                X                     X
[ Input x2 ] --/ \-- [ Hidden h2 ] --/ \-- [ Output y2 ]
      |                  |                  |
   (d-dim)           (k-dim)             (classes)
\end{verbatim}

\section{关键代码片段 (PyTorch实现)}
\begin{lstlisting}
import torch.nn as nn

class MLP(nn.Module):
    def __init__(self, input_size, hidden_size, num_classes):
        super(MLP, self).__init__()
        # 定义网络层序列
        self.network = nn.Sequential(
            # 第一层全连接 (Linear / Dense layer)
            nn.Linear(input_size, hidden_size),
            nn.ReLU(),
            # 第二层全连接
            nn.Linear(hidden_size, hidden_size),
            nn.ReLU(),
            # 输出层
            nn.Linear(hidden_size, num_classes)
        )

    def forward(self, x):
        # x shape: (batch_size, input_size)
        return self.network(x)
\end{lstlisting}

\section{深度面试题与解答}
\subsection{基础概念}
\textbf{问：为什么加入隐藏层和激活函数后，模型就能解决非线性问题？}\\
答：如果没有隐藏层的非线性激活，多层全连接网络在数学上等价于单层矩阵相乘，即 $W_2(W_1 x) = (W_2 W_1)x$，本质仍是线性模型。激活函数为每一层引入了扭曲空间的能力，多层叠加后可以构建极复杂的非线性决策边界。

\subsection{进阶原理}
\textbf{问：MLP 的参数量如何计算？}\\
答：设输入维度为 $D_i$，输出维度为 $D_o$。一层全连接层的参数量为：$D_i \times D_o (\text{Weights}) + D_o (\text{Biases})$。

\subsection{工程实战}
\textbf{问：MLP 在处理图像、文本时有什么局限性？}\\
答：1. \textbf{参数冗余}：全连接会忽略数据的空间（图像）或时序（文本）结构，导致参数量随输入尺寸爆炸式增长。2. \textbf{缺乏不变性}：无法利用图像的平移不变性，物体在图中移动一个像素，对于 MLP 来说就是完全不同的输入。

\subsection{坑点分析}
\textbf{问：如果隐藏层神经元非常多，会导致什么问题？}\\
答：容易产生\textbf{过拟合}（Overfitting）。模型会不仅学习到数据的真实规律，还会记住训练集的噪声。解决方法包括使用 Dropout、L2 正则化或减小隐藏层维度。

\section{MLP 与其他模型对比表}

\begin{table}[h]
\centering
\begin{tabularx}{\textwidth}{|l|X|}
\hline
\textbf{模型} & \textbf{核心特征对比} \\
\hline
\textbf{感知机} & 单层，无隐藏层，只能解决线性可分问题。 \\
\hline
\textbf{MLP} & 多层，万能近似器，但参数量大，不适合复杂结构数据。 \\
\hline
\textbf{CNN} & 局部连接、权重共享，擅长处理二维网格结构（图像）。 \\
\hline
\textbf{Transformer} & 基于注意力机制，擅长处理长程依赖和海量数据。 \\
\hline
\end{tabularx}
\end{table}

\section{参考文献与拓展}
\begin{enumerate}
    \item \textbf{经典文献}: Rumelhart, D. E., et al. (1986). \textit{Learning representations by back-propagating errors}. Nature.
    \item \textbf{深度学习圣经}: \textit{Deep Learning} (Ian Goodfellow, Chapter 6: Deep Feedforward Networks).
    \item \textbf{在线课程}: Andrew Ng's Neural Networks and Deep Learning (Coursera).
\end{enumerate}

\chapter{Positional\_encoding}


\section{核心背景：置换不变性}
\textbf{问题定义}：Self-Attention 机制本质上是加权和操作，对序列的顺序不敏感（置换不变性）。如果直接将词嵌入丢进 Transformer，模型将无法区分 "I love coding" 和 "coding love I"。因此必须显式地向模型注入位置信息。

\section{绝对位置编码 (Absolute PE)}
通过将位置向量与词嵌入向量相加（Addition）来实现。

\subsection{正余弦固定编码 (Sinusoidal)}
Vaswani 等人 (2017) 提出利用不同频率的正余弦函数：
$$ PE(pos, 2i) = \sin\left(\frac{pos}{10000^{2i/d_{model}}}\right), \quad PE(pos, 2i+1) = \cos\left(\frac{pos}{10000^{2i/d_{model}}}\right) $$

\textbf{设计动机：为什么选正余弦？}
\begin{itemize}
    \item \textbf{相对距离建模}：根据三角函数的和差化积公式，对于任何固定的偏移 $k$，$PE_{pos+k}$ 可以表示为 $PE_{pos}$ 的线性函数。这使得模型能更容易学习到词与词之间的固定的相对距离。
    \item \textbf{多尺度表达}：不同维度 $i$ 对应不同的波长，从 $2\pi$ 到 $10000 \cdot 2\pi$。这就像用时钟的秒针、分针、时针来共同表示一个精确的时间点，使得模型能够同时捕捉短程和长程的依赖。
\end{itemize}

\textbf{注入机制：它是如何生效的？}
当我们将词嵌入 $E$ 与位置编码 $P$ 相加进行 Self-Attention 计算时，点积项拆解为：
$$ (E_i + P_i)(E_j + P_j)^\top = \underbrace{E_i E_j^\top}_{内容-内容} + \underbrace{E_i P_j^\top}_{内容-位置} + \underbrace{P_i E_j^\top}_{位置-内容} + \underbrace{P_i P_j^\top}_{位置-位置} $$
这意味着位置信息参与了注意力权重的分配，使得 Query 能够根据“我在第 3 位，我要找第 1 位的东西”这种逻辑进行匹配。

\subsection{可学习编码 (Learned Embedding)}
直接为每个位置初始化一个可训练的向量 $W_{pos}$。BERT 使用此方法。
\textbf{局限}：训练时见过的最大长度是 $L$，推理时如果遇到 $L+1$ 处的位置索引，模型将无法处理（外推性差）。

\section{旋转位置嵌入 (RoPE)}
\textbf{现代 LLM (LLaMA, GLM) 的标配}。其核心是通过在复数空间进行旋转来编码位置。

\textbf{设计动机：相对位置的一致性}
我们希望找到一个函数 $f(x, m)$，使得两个向量点积的结果只依赖于它们的相对距离 $m-n$：
$$ \langle f(q, m), f(k, n) \rangle = g(q, k, m-n) $$
\textbf{数学实现}：对第 $m$ 个位置的向量在第 $i$ 对维度上乘以旋转矩阵：
$$ f(q, m) = \begin{pmatrix} \cos m\theta & -\sin m\theta \\ \sin m\theta & \cos m\theta \end{pmatrix} \begin{pmatrix} q_{2i} \\ q_{2i+1} \end{pmatrix} $$
\textbf{注入机制：旋转即位置}
由于旋转变换是正交变换，不改变向量的模长，但改变了两个向量之间的夹角。词与词之间的距离越远，旋转的角度差异越大，点积之后的结果就体现出了相对距离。

\section{线性偏置注意力 (ALiBi)}
Press 等人 (2021) 提出不再修改 Embedding，而是在注意力分数矩阵上直接叠加偏置。

\textbf{设计动机：极致的外推性}
$$ Attention(q_i, k_j) = \text{Softmax}(q_i k_j^T - m \cdot |i - j|) $$
\textbf{注入机制}：对于距离越远的词，人为地施加一个惩罚项。这种设计比模型自己去“学”长程规律要稳健得多。即使推理长度是训练长度的 10 倍，模型依然能保持基本性能，因为它不再依赖于绝对的位置索引。

\section{总结对比：如何选择？}
\begin{table}[h!]
\centering
\begin{tabularx}{\textwidth}{|l|X|X|}
\hline
\textbf{方案} & \textbf{注入方式} & \textbf{核心能力} \\
\hline
Sinusoidal & 加法融合 (Additive) & 相对距离建模能力。 \\
\hline
RoPE & 乘法旋转 (Multiplicative) & \textbf{外推性强}，保留语义相关性。 \\
\hline
ALiBi & 偏置干预 (Bias) & \textbf{外推性极强}，内存占用低。 \\
\hline
\end{tabularx}
\end{table}

\section{代码片段 (Sinusoidal 构造逻辑)}
\begin{lstlisting}
import torch
import math

def get_sinusoidal_pe(seq_len, d_model):
    pe = torch.zeros(seq_len, d_model)
    position = torch.arange(0, seq_len, dtype=torch.float).unsqueeze(1)
    div_term = torch.exp(torch.arange(0, d_model, 2).float() * (-math.log(10000.0) / d_model))
    # 为何这样写？e^(-i * log(10000)/d) = 10000^(-i/d)
    pe[:, 0::2] = torch.sin(position * div_term)
    pe[:, 1::2] = torch.cos(position * div_term)
    return pe
\end{lstlisting}

\section{深度面试题}
\textbf{问：为什么 Transformer 论文中说加法融合（$E+P$）不会破坏语义？}\\
答：因为词向量空间是非稠密的超高维空间。相加可以理解为在词向量中注入了一个“位置指向”的偏置，由于维数极高，这种加法操作在大数定律下类似于在词向量中保留了一个独立的通道来存储位置信息，后续的权重矩阵 $W^Q, W^K$ 在学习过程中会自动解析这两部分内容。

\chapter{Rnn}


\section{核心定义}
\textbf{中文定义}：循环神经网络（Recurrent Neural Network, RNN）是一类用于处理序列数据（如文本、时间序列、音频）的神经网络。与传统前馈网络不同，RNN 在层间引入了反馈循环，使得模型能够保留并利用历史信息。

\textbf{英文定义}：Recurrent Neural Networks (RNNs) are a class of neural networks designed for processing sequential data. Unlike feedforward networks, RNNs feature feedback loops, allowing them to maintain an internal state (memory) to process sequences of volatile length.

\textbf{核心价值}：
RNN 的核心在于\textbf{参数共享}。在序列的每一个时间步，模型都使用相同的权重矩阵来处理输入，这使得模型能够处理任意长度的序列。

\section{核心原理与常用变体}

\subsection{1. Vanilla RNN (基础循环网络)}
在每个时间步 $t$，隐藏状态 $h_t$ 的更新公式为：
$$ h_t = \sigma(W_{hh} h_{t-1} + W_{xh} x_t + b_h) $$
其中 $x_t$ 是当前输入，$h_{t-1}$ 是上一时刻的记忆。
\textbf{致命缺陷}：难以处理长距离依赖。在反向传播时，长期梯度会涉及多次权重矩阵相乘，导致\textbf{梯度消失}或\textbf{梯度爆炸}。

\subsection{2. LSTM (Long Short-Term Memory)}
通过引入\textbf{门控机制}（Gate Mechanism）解决长期依赖问题。
\begin{enumerate}
    \item \textbf{遗忘门} ($f_t$)：决定丢弃多少旧记忆。
    \item \textbf{输入门} ($i_t$)：决定存入多少新信息。
    \item \textbf{输出门} ($o_t$)：决定输出多少当前状态。
    \item \textbf{细胞状态} ($C_t$)：贯穿所有时间步的“主干线”，梯度可在此平滑流动。
\end{enumerate}

\subsection{3. GRU (Gated Recurrent Unit)}
LSTM 的简化版本，将遗忘门和输入门合并为\textbf{更新门} ($z_t$)，并将细胞状态与隐藏状态合并。计算开销更小，在中小规模数据集上表现常优于 LSTM。

\section{典型架构示意图}
RNN 的展开（Unrolling）结构如下：
\begin{verbatim}
      [ h_t-1 ]      [ h_t ]      [ h_t+1 ]
          ^              ^            ^
          |              |            |
(W_hh) -> [ Cell ] ----> [ Cell ] --> [ Cell ] -> (W_hh)
          ^              ^            ^
          |              |            |
      [ x_t-1 ]      [ x_t ]      [ x_t+1 ]
\end{verbatim}

\section{数学推导：BPTT 算法}
RNN 的训练算法称为\textbf{随时间反向传播（BPTT）}。总损失 $L$ 是各时刻损失之和。对于权重 $W_{hh}$ 的梯度：
$$ \frac{\partial L}{\partial W_{hh}} = \sum_{t=1}^T \frac{\partial L_t}{\partial W_{hh}} $$
在计算 $\frac{\partial L_t}{\partial W_{hh}}$ 时，需要追溯到 $t, t-1, \dots, 1$ 时刻的所有状态，梯度项包含 $\prod_{k=1}^t \frac{\partial h_k}{\partial h_{k-1}}$。若权重矩阵特征值不等于 1，连乘会导致数值剧变。

\section{关键代码片段 (PyTorch实现)}
\begin{lstlisting}
import torch
import torch.nn as nn

class RNNModel(nn.Module):
    def __init__(self, input_size, hidden_size, output_size):
        super().__init__()
        # input_size: 词向量维度, hidden_size: 记忆单元维度
        self.rnn = nn.RNN(input_size, hidden_size, batch_first=True)
        # 也可以替换为 nn.LSTM 或 nn.GRU
        # self.rnn = nn.LSTM(input_size, hidden_size, batch_first=True)
        
        self.fc = nn.Linear(hidden_size, output_size)

    def forward(self, x):
        # x shape: (batch, seq_len, input_size)
        # out: 包含所有时间步的隐藏状态 (batch, seq_len, hidden_size)
        # hn: 最后一个时间步的隐藏状态 (1, batch, hidden_size)
        out, hn = self.rnn(x)
        
        # 通常取最后一个时间步进行分类
        last_out = out[:, -1, :]
        return self.fc(last_out)
\end{lstlisting}

\section{深度面试题与解答}
\subsection{基础概念}
\textbf{问：RNN 为什么比全连接网络更适合处理文本？}\\
答：文本具有时序相关性（前后的词有逻辑联系）。RNN 的循环结构允许信息在时间步之间传递，实现了“记忆”功能。此外，文本长度是不固定的，RNN 的参数共享机制允许它处理任意长度的序列输入。

\subsection{进阶原理}
\textbf{问：LSTM 为什么能缓解梯度消失？}\\
答：因为 LSTM 引入了“细胞状态”路径，$C_t = f_t \odot C_{t-1} + i_t \odot \tilde{C}_t$。在反向传播时，梯度在 $C_t$ 通道上主要是相加关系，而不是连乘关系。即便某些门关闭，梯度依然可以通过遗忘门路径较好地传回较远的时间步。

\subsection{工程实战}
\textbf{问：双向 RNN (Bi-RNN) 的原理是什么？}\\
答：同时运行两个 RNN，一个从前向后（正向），一个从后向前（反向）。在每个时刻 $t$，将两者的隐藏状态拼接。这使得模型在预测当前词时，既能看到上文信息，也能看到下文信息。

\subsection{坑点分析}
\textbf{问：RNN 训练时最常见的数值问题是什么？如何解决？}\\
答：\textbf{梯度爆炸}。表现为 Loss 突然变成 NaN。标准解决方法是使用\textbf{梯度裁剪（Gradient Clipping）}，即如果梯度的范数超过某个阈值，则对其进行等比例缩放，防止更新步长过大。

\section{RNN vs Transformer}

\begin{table}[h]
\centering
\begin{tabularx}{\textwidth}{|l|X|X|}
\hline
\textbf{特性} & \textbf{RNN (LSTM/GRU)} & \textbf{Transformer} \\
\hline
计算模式 & 串行（逐个单词处理） & 并行（一次性处理全句） \\
\hline
长程依赖 & 依赖记忆门（有上限） & 全局自注意力（无损） \\
\hline
硬件效率 & 较低（GPU 利用率受限） & 极高 \\
\hline
位置信息 & 隐式（结构决定顺序） & 显式（需加入位置编码） \\
\hline
\end{tabularx}
\end{table}

\section{参考文献与拓展}
\begin{enumerate}
    \item \textbf{LSTM 奠基论文}: Hochreiter, S., & Schmidhuber, J. (1997). \textit{Long short-term memory}. Neural Computation.
    \item \textbf{GRU 论文}: Cho, K., et al. (2014). \textit{Learning Phrase Representations using RNN Encoder-Decoder}.
    \item \textbf{深入浅出}: \textit{Understanding LSTMs} (Colah's blog, 最经典的博文)。
\end{enumerate}

\chapter{Transformer}


\section{核心定义}
\textbf{中文定义}：Transformer 是一种基于自注意力机制（Self-Attention Mechanism）的深度学习架构，摒弃了传统的循环神经网络（RNN）和卷积神经网络（CNN），通过全局注意力机制捕获序列数据中的长距离依赖关系。

\textbf{英文定义}：Transformer is a deep learning architecture based entirely on self-attention mechanisms, dispensing with traditional recurrence and convolutions entirely to capture long-range dependencies in sequential data.

\textbf{提出时间与作者}：2017年，由 Google 团队（Ashish Vaswani 等人）提出。

\textbf{原始关键论文}：\textit{Attention Is All You Need} (NeurIPS 2017)

\textbf{技术体系定位及历史演进}：
Transformer 最初作为自然语言处理（NLP）领域机器翻译任务的 Seq2Seq 模型提出，旨在解决 RNN 难以并行计算以及长跨度依赖（Long-term Dependency）容易遗忘的问题。由于其卓越的特征提取能力和极强的并行扩展性，Transformer 迅速取代了 LSTM/GRU，成为 NLP 的事实标准（催生了 BERT、GPT 等大模型）。随后，其应用领域扩展至计算机视觉（Vision Transformer, ViT）、多模态学习以及强化学习等各个 AI 领域，成为现代深度学习的基础架构。

\section{核心原理与底层逻辑}
Transformer 的核心在于\textbf{多头自注意力（Multi-Head Self-Attention）}机制，加上前馈神经网络（Feed-Forward Networks）、残差连接（Residual Connection）和层归一化（Layer Normalization）。

\subsection{数据流向简述}
假定输入序列长度为 $L$，嵌入维度为 $D$，批量大小为 $B$。
\begin{enumerate}
    \item \textbf{输入张量}：原始词元进入 Embedding 层，转化为形状为 $(B, L, D)$ 的张量，并加上相同形状的 \textbf{位置编码（Positional Encoding）}，以注入序列的顺序信息。
    \item \textbf{自注意力计算}：张量通过线性变换分别投影出 $Q$ (Query), $K$ (Key), $V$ (Value)，形状均为 $(B, L, D)$（多头情况下会被拆分成多份）。注意力机制计算词元之间的相关性得分并加权求和，输出形状依旧是 $(B, L, D)$。
    \item \textbf{前馈网络与输出}：经过残差连接与 LayerNorm 后，送入两层全连接网络（内部维度通常放大 4 倍），再经残差与 LayerNorm 最终输出形状为 $(B, L, D)$ 的张量。
\end{enumerate}

\subsection{算法执行步序}
\begin{enumerate}
    \item 将输入 Token 映射为词向量 $X$。
    \item 生成位置编码 $PE$，并计算 $X \leftarrow X + PE$。
    \item 将 $X$ 投影为三个矩阵：查询（$Q$）、键（$K$）、值（$V$）。
    \item 计算缩放点积注意力（Scaled Dot-Product Attention）：$Q$ 与 $K^T$ 相乘，除以放缩因子，经过 Softmax 得到注意力权重。
    \item 注意力权重与 $V$ 相乘，得到加权上下文向量。
    \item 加入残差连接并进行 LayerNorm：$\text{Output}_1 = \text{LayerNorm}(X + \text{Attention}(X))$。
    \item 输入到全连接前馈网络（FFN），再次加入残差连接并进行 LayerNorm：$\text{Output}_2 = \text{LayerNorm}(\text{Output}_1 + \text{FFN}(\text{Output}_1))$。
\end{enumerate}

\section{严谨数学公式与推导}
Transformer 的最核心公式是缩放点积注意力（Scaled Dot-Product Attention）：
$$
\text{Attention}(Q,K,V) = \text{softmax}\left(\frac{QK^T}{\sqrt{d_k}}\right)V
$$
\textbf{符号说明}：
\begin{itemize}
    \item $Q \in \mathbb{R}^{L \times d_k}$：Query 矩阵，表示当前词元的查询请求。
    \item $K \in \mathbb{R}^{L \times d_k}$：Key 矩阵，表示序列中所有词元的被查询标识。
    \item $QK^T \in \mathbb{R}^{L \times L}$：两者相乘得到注意力分数矩阵，表示序列中第 $i$ 个词元对第 $j$ 个词元的注意力大小。
    \item $\sqrt{d_k}$：缩放因子（Scaling Factor），$d_k$ 是 $K$ 的维度。作用是防止点积结果过大，导致 Softmax 函数进入梯度极小的饱和区。
    \item $V \in \mathbb{R}^{L \times d_v}$：Value 矩阵，表示词元的实际内容表示。
    \item \text{softmax}：对每一行（每个词元的注意力分布）进行归一化，使其和为 1。
\end{itemize}

多头注意力（Multi-Head Attention）则是将 $Q, K, V$ 映射到多个不同的线性子空间中分别计算注意力，再将结果拼接融合：
$$
\text{MultiHead}(Q,K,V) = \text{Concat}(\text{head}_1, \dots, \text{head}_h) W^O
$$
其中，$\text{head}_i = \text{Attention}(Q W_i^Q, K W_i^K, V W_i^V)$。

\section{模型结构示意图 (Full Architecture)}

\begin{figure}[h!]
\centering
\resizebox{0.88\textwidth}{!}{
\begin{tikzpicture}[node distance=1.8cm, auto]

    % --- Encoder ---
    \node [emb] (input_emb) {Input Embedding};
    \node [pos] (pos_enc_in) [right=0.2cm of input_emb] {+};
    \node [below=0.1cm of pos_enc_in, font=\tiny] {Positional Encoding};
    
    \node [attn, above=of input_emb] (multi_head_enc) {Multi-Head Self-Attention};
    \node [norm, above=0.6cm of multi_head_enc] (norm1_enc) {Add \& Norm};
    \node [ffn, above=of norm1_enc] (ffn_enc) {Feed Forward};
    \node [norm, above=0.6cm of ffn_enc] (norm2_enc) {Add \& Norm};
    
    \draw [arrow] (input_emb) -- (multi_head_enc);
    \draw [arrow] (multi_head_enc) -- (norm1_enc);
    \draw [arrow] (norm1_enc) -- (ffn_enc);
    \draw [arrow] (ffn_enc) -- (norm2_enc);
    
    % Residuals Encoder
    \draw [arrow] ($(input_emb.north)$) -- ++(-1.8,0) |- (norm1_enc.west);
    \draw [arrow] ($(norm1_enc.north)$) -- ++(-1.8,0) |- (norm2_enc.west);
    
    \node [left=1.5cm of norm2_enc, font=\small\bfseries] {Encoder ($\times N$)};

    % --- Decoder ---
    \node [emb, right=4.0cm of input_emb] (output_emb) {Output Embedding\\(Shifted Right)};
    \node [pos] (pos_enc_out) [right=0.2cm of output_emb] {+};
    
    \node [attn, above=of output_emb] (masked_attn) {Masked Multi-Head\\Self-Attention};
    \node [norm, above=0.6cm of masked_attn] (norm1_dec) {Add \& Norm};
    
    \node [attn, above=of norm1_dec] (cross_attn) {Multi-Head\\Cross-Attention};
    \node [norm, above=0.6cm of cross_attn] (norm2_dec) {Add \& Norm};
    
    \node [ffn, above=of norm2_dec] (ffn_dec) {Feed Forward};
    \node [norm, above=0.6cm of ffn_dec] (norm3_dec) {Add \& Norm};
    
    \node [layer, fill=purple!10, above=of norm3_dec] (linear) {Linear \& Softmax};
    
    \draw [arrow] (output_emb) -- (masked_attn);
    \draw [arrow] (masked_attn) -- (norm1_dec);
    \draw [arrow] (norm1_dec) -- (cross_attn);
    \draw [arrow] (cross_attn) -- (norm2_dec);
    \draw [arrow] (norm2_dec) -- (ffn_dec);
    \draw [arrow] (ffn_dec) -- (norm3_dec);
    \draw [arrow] (norm3_dec) -- (linear);
    
    % Residuals Decoder
    \draw [arrow] ($(output_emb.north)$) -- ++(1.8,0) |- (norm1_dec.east);
    \draw [arrow] ($(norm1_dec.north)$) -- ++(1.8,0) |- (norm2_dec.east);
    \draw [arrow] ($(norm2_dec.north)$) -- ++(1.8,0) |- (norm3_dec.east);
    
    % Cross-Attention Connection
    \path (norm2_enc.east) -- (cross_attn.west) coordinate[midway] (mid_point);
    \draw [arrow, dashed, ultra thick, blue!60] (norm2_enc.east) -- (mid_point |- norm2_enc.east) 
          node[above, font=\tiny, fill=white, inner sep=1pt] {Encoder Output (K, V)} 
          -| (cross_attn.west);
    
    \node [right=1.0cm of norm3_dec, font=\small\bfseries] {Decoder ($\times N$)};

\end{tikzpicture}}
\caption{Transformer Full Architecture (Attention Is All You Need)}
\end{figure}

\section{关键代码片段 (PyTorch实现)}
以下为 Transformer 中最核心的多头自注意力的简化实现（便于展示原理，省略了多头拆分的重塑操作细节）：

\begin{lstlisting}
import torch
import torch.nn as nn
import torch.nn.functional as F
import math

class SelfAttention(nn.Module):
    def __init__(self, d_model, heads):
        super().__init__()
        self.d_model = d_model
        self.heads = heads
        self.d_k = d_model // heads
        
        # 定义计算 Q, K, V 的线性层
        self.q_linear = nn.Linear(d_model, d_model)
        self.k_linear = nn.Linear(d_model, d_model)
        self.v_linear = nn.Linear(d_model, d_model)
        
        # 最后的输出映射层
        self.out = nn.Linear(d_model, d_model)
        
    def forward(self, x, mask=None):
        # x shape: (B, L, d_model)
        B, L, _ = x.shape
        
        # 1. 线性映射并拆分为多头 (B, L, heads, d_k) -> 调换维度 (B, heads, L, d_k)
        # 为了代码简洁，用伪操作表示 reshape 和 transpose
        Q = self.q_linear(x).view(B, L, self.heads, self.d_k).transpose(1, 2)
        K = self.k_linear(x).view(B, L, self.heads, self.d_k).transpose(1, 2)
        V = self.v_linear(x).view(B, L, self.heads, self.d_k).transpose(1, 2)
        
        # 2. 计算 Scaled Dot-Product Attention: Q * K^T / sqrt(d_k)
        # attention shape: (B, heads, L, L)
        scores = torch.matmul(Q, K.transpose(-2, -1)) / math.sqrt(self.d_k)
        
        if mask is not None:
            scores = scores.masked_fill(mask == 0, -1e9)
            
        # 注意力权重归一化
        attn_weights = F.softmax(scores, dim=-1)
        
        # 3. 乘以 Value
        # output shape: (B, heads, L, d_k)
        context = torch.matmul(attn_weights, V)
        
        # 4. 拼接多头并线性映射恢复原始维度
        context = context.transpose(1, 2).contiguous().view(B, L, self.d_model)
        output = self.out(context) # (B, L, d_model)
        
        return output
\end{lstlisting}

\section{深度面试题与解答}
\subsection{基础概念}
\textbf{问：Transformer 为什么要加入位置编码及其运作方式？}\\
答：因为自注意力机制（Q、K、V相乘）计算两两之间的关系时不关注它们在序列中的物理距离，这是一种完全等距的“集合”操作。如果不加入位置编码，模型会将“I like ML”和“ML like I”视作完全一致。原论文采用了正弦和余弦函数的绝对位置编码，通过将位置编码直接与输入词嵌入相加（$dim$一致）来给每个位置打上标签。

\subsection{进阶原理}
\textbf{问：为什么点积注意力在经过 Softmax 之前要除以 $\sqrt{d_k}$？}\\
答：当维度 $d_k$ 很大时，相互独立且服从标准正态分布的 $Q$ 和 $K$ 的点积，均值虽为 0，但方差会增大至 $d_k$。巨大的方差导致分布呈现两极分化，某些元素值极大。此时传入 Softmax 函数，会导致大的部分接近 1、其余接近 0，即 Softmax 函数陷入了梯度极小的平滑区域。除以 $\sqrt{d_k}$ 相当于将点积结果的方差缩放回 1，使得梯度流更加稳定。

\subsection{工程实战}
\textbf{问：在 Transformer 中 LayerNorm 放置的位置有哪几种？各有什么优劣？}\\
答：主要分为 Post-LN（原论文，在残差相加后执行 LN）和 Pre-LN（目前的绝对主流如 GPT-3，在 Attention 和 FFN 之前执行 LN）。\textbf{Post-LN} 对权重的初始分布极其敏感，梯度容易在靠近输入端消失，训练极不稳定，常需要设置复杂的 Warmup 学习率调度机制；\textbf{Pre-LN} 的梯度流动从后到前是一致的，不管层数多深，无需 Warmup 也可以直接稳定收敛。

\subsection{坑点分析}
\textbf{问：Transformer 处理长文本时面临的最大性能瓶颈是什么？如何解决？}\\
答：瓶颈在于计算自注意力的复杂度是 $O(L^2 \times D)$，空间复杂度（存储 Attention Map）也是 $O(L^2)$。当需要处理数十万 token 时显存会瞬间爆炸。解决方式有：
1. 稀疏注意力（如 Longformer 仅计算局部窗）。
2. 线性关联矩阵（如 Linear Transformer 重新排列向量乘积顺序）。
3. 状态复用或内存压缩机制（如 Transformer-XL）。
4. (近期绝对主流) 状态空间模型（Mamba）或是 FlashAttention 高效 IO 并发算法。

\section{优缺点与适用场景}

\begin{table}[h]
\centering
\begin{tabular}{|p{0.2\textwidth}|p{0.75\textwidth}|}
\hline
\textbf{维度} & \textbf{分析} \\
\hline
\textbf{优势} & 1. 全局视野：一步直接建立任意长距离依赖关系。\\
& 2. 高度并行化：训练时不需要像 RNN 那样等待前序步骤完成，硬件利用率极高。\\
& 3. 泛化与扩展性：模型容量扩展到千亿参数也未见明显天花板（Scaling Law 表观有效）。\\
\hline
\textbf{局限} & 1. 计算复杂度过高：序列长度的二次方复杂度阻碍了极致长文应用。\\
& 2. 缺乏归纳偏置（Inductive Bias）：相比CNN对图像平移的不变性假设，Transformer 较难在小数据集上收敛，极为依赖海量数据的暴力压制。\\
\hline
\textbf{最佳适用场景} & 海量数据下的大模型预训练（NLP 的语言模型、代码生成）、机器翻译、高质量文本生成；近期的图像生成的 DiT 架构。\\
\hline
\textbf{绝对慎用场景} & 1. 训练数据极少且无预训练支持的孤立小任务。\\
& 2. MCU（单片机）等边缘侧算力极度受限设备的实时推理。\\
\hline
\end{tabular}
\end{table}

\section{参考文献与拓展}
\begin{enumerate}
    \item \textbf{核心经典论文}: Vaswani, A., et al. (2017). \textit{Attention is all you need}. NeurIPS. 
    \item \textbf{代码实现参考}: \texttt{torch.nn.MultiheadAttention} 和 \texttt{torch.nn.TransformerEncoderLayer} 官方源码。
    \item \textbf{必读进阶}: \textit{The Illustrated Transformer}（Jay Alammar，目前最生动的图文详解）。
    \item \textbf{变体探索}: FlashAttention (Dao et al., 2022) 解决硬件 IO 通信瓶颈的终极方案。
\end{enumerate}

\chapter{Vae}


\section{核心定义}
\textbf{中文定义}：变分自编码器（Variational Autoencoder, VAE）是一种生成模型，它将自编码器（AE）与变分推断相结合。不同于普通 AE 将输入映射为固定的隐向量，VAE 将输入映射为一个概率分布（通常是正态分布），从而使得隐空间具有连续性和可采样性。

\textbf{英文定义}：Variational Autoencoders (VAEs) are generative models that combine autoencoders with variational inference. Unlike standard AEs that map inputs to discrete latent points, VAEs map inputs to a probability distribution in the latent space, ensuring continuity and allowing for sampling.

\textbf{提出时间与作者}：2013年，由 Diederik P. Kingma 和 Max Welling 提出。

\textbf{原始关键论文}：\textit{Auto-Encoding Variational Bayes} (ICLR 2014)

\section{核心原理：潜在空间参数化}

\subsection{1. 结构组成}
\begin{itemize}
    \item \textbf{编码器 (Encoder)}：学习映射 $q_\phi(z|x)$，输出均值 $\mu$ 和方差 $\sigma$（通常输出 $\log \sigma^2$ 以保证数值稳定）。
    \item \textbf{解码器 (Decoder)}：学习映射 $p_\theta(x|z)$，从隐变量 $z$ 中重构输入 $\hat{x}$。
\end{itemize}

\subsection{2. 重参数化技巧 (Reparameterization Trick)}
由于从分布 $\mathcal{N}(\mu, \sigma^2)$ 中直接采样是一个不可导的操作，无法进行反向传播。VAE 引入了重参数化：
$$ z = \mu + \sigma \odot \epsilon, \quad \epsilon \sim \mathcal{N}(0, \mathbf{I}) $$
这样，随机性被转移到了外部变量 $\epsilon$ 上，而梯度可以顺利通过 $\mu$ 和 $\sigma$ 回传。

\section{损失函数：ELBO}
VAE 的训练目标是最大化证据下界（Evidence Lower Bound, ELBO），其损失函数由两部分组成：
$$ \mathcal{L} = \mathcal{L}_{reconstruction} + \mathcal{L}_{KL} $$
\begin{enumerate}
    \item \textbf{重构损失}：衡量输出与输入的相似度（通常使用 MSE 或交叉熵）。
    \item \textbf{KL 散度损失}：约束隐分布 $q(z|x)$ 逼近标准正态分布 $\mathcal{N}(0, \mathbf{I})$。公式为：
    $$ D_{KL}( \mathcal{N}(\mu, \sigma^2) || \mathcal{N}(0, \mathbf{I}) ) = -\frac{1}{2} \sum (1 + \log(\sigma^2) - \mu^2 - \sigma^2) $$
\end{enumerate}

\section{模型结构示意图}
\begin{verbatim}
[输入 x] --> [Encoder] --> [mu, log_var] --(采样 z = mu + exp(0.5*log_var)*eps)-->
                                  |
                                  v
                            [Decoder] --> [重构输出 x_hat]
\end{verbatim}

\section{关键代码片段 (PyTorch实现)}
\begin{lstlisting}
import torch
import torch.nn as nn

class VAE(nn.Module):
    def __init__(self, input_dim, latent_dim):
        super(VAE, self).__init__()
        # Encoder
        self.fc_e = nn.Linear(input_dim, 400)
        self.fc_mu = nn.Linear(400, latent_dim)
        self.fc_logvar = nn.Linear(400, latent_dim)
        # Decoder
        self.fc_d1 = nn.Linear(latent_dim, 400)
        self.fc_d2 = nn.Linear(400, input_dim)

    def encode(self, x):
        h = torch.relu(self.fc_e(x))
        return self.fc_mu(h), self.fc_logvar(h)

    def reparameterize(self, mu, logvar):
        std = torch.exp(0.5 * logvar)
        eps = torch.randn_like(std)
        return mu + eps * std

    def decode(self, z):
        h = torch.relu(self.fc_d1(z))
        return torch.sigmoid(self.fc_d2(h))

    def forward(self, x):
        mu, logvar = self.encode(x.view(-1, 784))
        z = self.reparameterize(mu, logvar)
        return self.decode(z), mu, logvar
\end{lstlisting}

\section{深度面试题与解答}
\subsection{基础概念}
\textbf{问：VAE 和普通 Autoencoder (AE) 的核心区别是什么？}\\
答：普通 AE 将输入映射为一个特定的点，隐空间往往是不连续的，无法直接用于生成；VAE 将输入映射为一个区域（分布），通过 KL 散度约束，确保了隐空间的数学完备性，使得我们可以从隐空间随机采样出有意义的新样本。

\subsection{进阶原理}
\textbf{问：KL 散度项在 VAE 中起什么作用？}\\
答：它起到正则化作用。如果没有 KL 项，编码器为了减小重构误差，会将 $\sigma$ 缩减到极小，使 VAE 退化为普通 AE。KL 项迫使隐分布向标准正态分布靠拢，从而在隐空间中保留了一定的“容错性”和“平滑性”。

\subsection{工程实战}
\textbf{问：VAE 生成的图像为什么通常比较模糊？}\\
答：因为 VAE 的损失函数假设数据分量之间是独立的（如 MSE 损失），且它在隐空间建模时采用的是连续的高斯分布，倾向于对数据分布进行平均化处理。相比之下，GAN 由于采用对抗判别，能捕捉更尖锐的边缘和细节。

\subsection{坑点分析}
\textbf{问：什么是后验崩溃 (Posterior Collapse)？}\\
答：在训练强大的解码器（如带有自回归性质的解码器）时，模型可能会发现直接忽略隐变量 $z$ 也能取得不错的重构效果，导致编码器学到的分布与先验完全一致（KL 散度变为 0）。解决方法包括 KL 退火、使用更简单的解码器等。

\section{AE vs VAE 对比表}

\begin{table}[h]
\centering
\begin{tabularx}{\textwidth}{|l|X|X|}
\hline
\textbf{维度} & \textbf{Autoencoder (AE)} & \textbf{Variational AE (VAE)} \\
\hline
隐空间性质 & 离散、不规则 & 连续、正态分布化 \\
\hline
核心目标 & 纯粹的降维与重构 & 学习数据生成的概率分布 \\
\hline
生成能力 & 弱（点对点映射） & 强（可从隐空间采样） \\
\hline
数学基础 & 几何重构误差 & 贝叶斯统计与变分推断 \\
\hline
\end{tabularx}
\end{table}

\section{参考文献与拓展}
\begin{enumerate}
    \item \textbf{奠基论文}: Kingma, D. P., & Welling, M. (2013). \textit{Auto-encoding variational bayes}.
    \item \textbf{进阶阅读}: \textit{Tutorial on Variational Autoencoders} (Carl Doersch, 2016)。
    \item \textbf{应用案例}: Conditional VAE (CVAE) 实现受控图像生成。
\end{enumerate}

\end{document}